% Source: https://artofproblemsolving.com/community/c3h1637714_random_problemo
% Stable identity: aops-topic-1637714-q01
\textbf{First circle.} A circle inscribed in a square of side $12$ has diameter $12$, so radius $6$.

\textbf{Second square.} A square inscribed in a circle of radius $6$ has diagonal $12$, so its side is $\dfrac{12}{\sqrt2}=6\sqrt2$.

\textbf{Second circle.} Its diameter is the side of the second square, so its radius is $3\sqrt2$.

\textbf{Triangle.} An equilateral triangle inscribed in a circle of radius $R$ has its centre at the centre of the circle, and each vertex is at distance $R=\dfrac{2}{3}\cdot\text{(height)}$ from it. So the height is $\dfrac{3}{2}R$ and the side is $a=\dfrac{2}{\sqrt3}\cdot\dfrac{3}{2}R=\sqrt3R$. With $R=3\sqrt2$, $a=3\sqrt6$ and
\[
\text{Area}=\dfrac{\sqrt3}{4}a^2=\dfrac{\sqrt3}{4}\cdot54=\boxed{\dfrac{27\sqrt3}{2}}.
\]

\begin{asy}
usepackage("amsmath");
size(7cm);
defaultpen(fontsize(10pt));
real s=6;real r=3*sqrt(2);
draw((-s,-s)--(s,-s)--(s,s)--(-s,s)--cycle,linewidth(0.9));
draw(circle((0,0),s));
draw((s,0)--(0,s)--(-s,0)--(0,-s)--cycle,linewidth(0.9));
draw(circle((0,0),r));
pair P=r*dir(90),Q=r*dir(210),R=r*dir(330);
filldraw(P--Q--R--cycle,paleblue,blue+linewidth(0.9));
draw((0,0)--(0,-s),red+dashed);label("$6$",(0,-3.1),E,red);
draw((0,0)--P,red+dashed);label("$3\sqrt2$",(0,1.3),W,red+fontsize(9pt));
dot((0,0));
label("$12$",(0,-s),S);
label("\mbox{}",(min(currentpicture,user=true)+max(currentpicture,user=true))/2,black); // keeps the text after the figure black when LaTeX typesets the labels inline
\end{asy}

(If instead the second circle were inscribed \emph{in} the triangle, the triangle would have inradius $3\sqrt2$, side $6\sqrt6$ and area $54\sqrt3$.)
